\chapter{Fazit}

Bei den Messungen mit Braggstreuung wurde die Kristallstruktur von NaCl analysiert und die Unbekannte Anode als Wolframanode identifiziert.
Bei der Röntgenfluoreszenspektroskopie wurden die Materialzusammensetzungen der unbekannten Proben gut bestimmt und es wurde nachvollziehbare, Handelsübliche Legierungen (Nickelroheisen, Messing und Bronze) gefunden, sowie die Massenverhältnisse der Bestandteile bestimmt.

Beim abschließenden Laue verfahren wurde die Messmethode nachvollzogen, wobei eine eindeutige Zurodnung nur schwer möglich war. Es wurden Netzebenabstände, Winkel und Wellenlängen identifiziert, welche weitesgehend Plausibel waren.
